% -*- coding: utf-8 -*-
% !TEX program = xelatex
\documentclass[plain,noanswer,medmath]{../../template/jnuexam}
\usepackage{../../template/kysx}

\begin{document}


\examtitle{1989}{5}


\loadproblems{1}{1989P1}%载入试卷一
\loadproblems{4}{1989P4}%载入试卷四


\makepart{填空题}{本题满分15分，每小题3分}


\useproblem{4}{1}{1}


\begin{problem}
某商品的需求量$Q$与价格$p$的函数关系为$Q=ap^b$，
其中$a$和$b$为常数，且$a\ne 0$，则需求量对价格$p$的弹性是 \fillin{}.
\end{problem}


\begin{problem}
行列式$\left| \begin{matrix}
   1 &  -1 &   1 & x-1 \\
   1 &  -1 & x+1 &  -1 \\
   1 & x-1 &   1 &  -1 \\
 x+1 &  -1 &   1 &  -1 \\
\end{matrix} \right|$= \fillin{}.
\end{problem}


\begin{problem}
设随机变量$X_1,X_2,X_3$相互独立，其中$X_1$在$[0,6]$上服从均匀分布，
$X_2$服从正态分布$N(0,2^2)$，$X_3$服从参数为$\lambda =3$的泊松分布，
记$Y=X_1-2X_2+3X_3$，则$DY=$ \fillin{}.
\end{problem}


\useproblem{4}{1}{4}


\makepart{选择题}{本题满分15分，每小题3分}


\useproblem{4}{2}{1}%【同试卷四第二(1)题】


\useproblem{4}{2}{2}%【同试卷四第二(2)题】


\useproblem{1}{2}{5}%【同试卷一第二(5)题】
%\useproblem{4}{2}{3}%【同试卷四第二(3)题】


\begin{problem}
设$n$元齐次线性方程组$AX=0$的系数矩阵$A$的秩为$r$，则$AX=0$有非零解的充分必要条件是 \pickout{}
\begin{abcd}
\item $r=n$．
\item $r<n$．
\item $r\ge n$．
\item $r>n$．
\end{abcd}
\end{problem}


\useproblem{4}{2}{5}%【同试卷四第二(5)题】


\makepart{计算题}{本题满分20分，每小题5分}

\begin{problem}
求极限$\lim_{x\to +\infty}\,(x+\e^x)^{\frac{1}{x}}$．
\end{problem}


\begin{problem}
已知$z=a^{\sqrt{x^2-y^2}}$，其中$a>0$, $a\ne 1$，求$\dz$．
\end{problem}


\begin{problem}
求不定积分$\int\frac{x+\ln(1-x)}{x^2}\dx$．
\end{problem}


\begin{problem}
求二重积分$\iint_D \frac{1-x^2-y^2}{1+x^2+y^2}\dx\dy$，
其中$D$是$x^2+y^2=1$,$x=0$和$y=0$所围成的区域在第Ⅰ象限部分．
\end{problem}


\makepart{}{本题满分6分}

\begin{problem}
已知某企业的总收入函数为$R=26x-2x^2-4x^3$，总成本函数为$C=8x+ x^2$，
其中$x$ 表示产品的产量，求利润函数、边际收入函数、边际成本函数、
以及企业获得最大利润时的产量和最大利润．
\end{problem}


\makepart{}{本题满分12分}

\begin{problem}
已知函数$y=\frac{2x^2}{(1-x)^2}$，试求其单调区间、极值点、
及图形的凹凸性、拐点和渐近线，并画出函数的图形．
\end{problem}


\makepart{}{本题满分5分}

\useproblem{4}{7}{0}%【同试卷四第七题】


\makepart{}{本题满分6分}%【同试卷四第八题】

\useproblem{4}{8}{0}


\makepart{}{本题满分5分}%【同试卷四第九题】

\useproblem{4}{9}{0}


\makepart{}{本题满分8分}
\begin{problem}
已知随机变量$X$和$Y$的联合分布为
$$\begin{array}{|c|cccccc|}
\hline
         (x,y) & (0,0) & (0,1) & (1,0) & (1,1) & (2,0) & (2,1) \\
\hline
  P\{X=x,Y=y\} &  0.10 &  0.15 &  0.25 &  0.20 &  0.15 & 0.15 \\
\hline
\end{array}$$
试求：\par
\begin{enumerate*}
\item $X$的概率分布；
\item $X+Y$的概率分布；
\item $Z=\sin\frac{\pi(X+Y)}{2}$的数学期望．
\end{enumerate*}
\end{problem}


\makepart{}{本题满分8分}
\begin{problem}
某仪器装有三只独立工作的同型号电子元件，其寿命（单位：小时）都服从同一指数分布，分布密度为
$$f(x)=\begin{cases}
 \tfrac{1}{600}\e^{-x/600}, & x>0, \\
 0, & x\le 0. \\
\end{cases}$$
试求：在仪器使用的最初$200$小时内，至少有一只电子元件损坏的概率$\alpha$．
\end{problem}


\end{document}
